% Copyright 2022 Thomas Ascher
% SPDX-License-Identifier: CC-BY-SA-4.0

\documentclass[a4paper,parskip=half]{scrartcl}

\usepackage[T1]{fontenc}
\usepackage{mathpazo}
\usepackage[naustrian]{babel}
\usepackage{csquotes}
\usepackage{booktabs}
\usepackage{graphicx}
\usepackage{chemformula}
\usepackage{icomma}
\usepackage{gensymb}
\usepackage{float}
\usepackage[section]{placeins}
\usepackage[style=apa,backend=biber]{biblatex}

\usepackage[hidelinks,pdfencoding=auto,
  pdfauthor={Thomas Ascher},
  pdfusetitle,
  pdfkeywords={Bier,Farbe,SRM,IBU,EBC,Morey}]{hyperref}
\usepackage{microtype}

\setkomafont{disposition}{\normalfont\bfseries}

\addto\extrasnaustrian{
\def\figureautorefname{Abb.}
\def\tableautorefname{Tab.}
\def\equationautorefname{Gl.}
}

\addto\captionsnaustrian{
\renewcommand{\figurename}{Abb.}
\renewcommand{\tablename}{Tab.}
}

\NewBibliographyString{gethesis}
\DefineBibliographyStrings{naustrian}{
  mathesis = {Masterarbeit},
  gethesis = {Diplomarbeit},
}

\newcommand{\MCU}{\mathit{MCU}}
\newcommand{\SRM}{\mathit{SRM}}
\newcommand{\umin}{\:[\textrm{min}]}
\newcommand{\uden}{\:[\text{g/cm³}]}
\newcommand{\uper}{\:[\text{\%}]}
\newcommand{\uli}{\:[\text{l}]}
\newcommand{\ume}{\:[\text{m}]}
\newcommand{\ucon}{\:[\text{mg/l}]}
\newcommand{\FKd}{F_{\mathit{d}}}

\title{Es wird bunt: Bierfarbe aus der Schüttung berechnen}
\author{Thomas Ascher <thomas.ascher@gmx.at>}
\date{\today, \href{http://creativecommons.org/licenses/by-sa/4.0/}{CC BY-SA 4.0}}

\addbibresource{bierfarbe.bib}

\begin{document}
\maketitle

\section*{Einleitung}

Bernsteinfarben, Schwarz, Rot oder Braun: Abseits der blonden Lager existiert das Produkt Bier innerhalb eines weitreichenden Farbspektrums. Neben einem nicht weltweit normierten Vokabular, wie es zum Beispiel \autoref{table:bacolor} zeigt, werden auch verschiedene Skalen verwendet, um diese Farben zu beschreiben. Großteils verantwortlich für dieses Farbenspiel ist die in einem Rezept eingesetzte Malzmischung, die Schüttung. Wie aus dieser die zu erwartende Bierfarbe näherungsweise berechnet werden kann, wird in den folgenden Abschnitten dieses Artikels demonstriert.

\begin{table}[H]
\centering
\begin{tabular}{lrr}
\toprule
\multicolumn{1}{c}{\textbf{Farbbeschreibung}} & \multicolumn{1}{c}{\textbf{SRM}} & \multicolumn{1}{c}{\textbf{EBC}} \\
\midrule
Sehr hell & 1–1,5 & 2–3  \\
Stroh & 2–3 & 4–6 \\
Blass & 4 & 8 \\
Gold & 5–6 & 10–12 \\
Helles Bernstein & 7 & 14 \\
Bernstein & 8 & 16 \\
Dunkles Bernstein & 9 & 18 \\
Kupfer/Granatrot & 10–12 & 20–24 \\
Hellbraun & 13–15 & 26–30 \\
Braun/Rotbraun/Kastanienbraun & 16–17 & 32–33 \\
Dunkelbraun & 18–24 & 35–47 \\
Sehr dunkel & 25–39 & 49–77 \\
Schwarz & 40+ & 79 \\
\bottomrule
\end{tabular}
\caption{Skalenwerte für Farbeindrücke \parencite{BA2021}}
\label{table:bacolor}
\end{table}

\begin{figure}[h]
\centering
\includegraphics[width=14cm]{colorscale.pdf}
\caption{Farbsimulation für Werte der \autoref{table:bacolor} (Ascher, 2022)}
\label{fig:mcucompare}
\end{figure}



\section*{Messung von Bier- und Malzfarbe}

Zur Bestimmung der Bierfarbe existieren mehrere Messskalen und -verfahren. Die Älteste davon noch gebräuchliche, ist die Lovibond Skala. Diese wurde vom Brauer Josep W. Lovibond 1885 zusammen mit einem Messgerät, dem Tintometer, entwickelt. Damit wird per visuellem Eindruck von einer Person eine aufbereitete Bierprobe mit verschieden gefärbten Glasscheiben verglichen, um einen Messwert zwischen 2 bis 70~°L (hell bis dunkel) zu ermitteln. Unterschiede in der menschlichen Wahrnehmung und verschiedene Umwelteinflüsse verfälschen das Messergebnis. \parencite{KrausWeyermann2021a}

Im Jahr 1951 hat die \href{https://www.asbcnet.org}{American Society of Brewing Chemists (ASBC)} eine neue Analysemethode (Standard Reference Method) auf Basis des Spektrophotometers zusammen mit der SRM Skala, die an die Lovibond Skala gekoppelt ist, ratifiziert. Seit 1996 verwendet die \href{https://europeanbreweryconvention.eu}{European Brewery Convention (EBC)} die EBC Skala, die sich ebenfalls auf die Messwerte eines Spektrophotometers bezieht. Bis auf die Skala wurde die grundlegende Messmethode von beiden Normungsgremien mittlerweile harmonisiert. Bei dieser misst ein Spektrophotometer die Lichtabsorption einer entgasten Bierprobe oder eine aus einem Malz im Kongressmaischeverfahren hergestellten Würzeprobe bei einer Wellenlänge von 430~nm durch eine 1~cm breite Quarzglasküvette. Der 12,7-fache Logarithmus der gemessenen Absorption entspricht dabei einem Wert der SRM Skala und der 25-fache Logarithmus der EBC Skala. Das menschliche Auge kann Farbunterschiede ab 80~EBC nicht mehr wahrnehmen. Auch Messgeräte generieren darüber keine zuverlässigen Messwerte mehr. Dunkle Bierproben sind daher vor der Messung zu verdünnen. \parencite{KrausWeyermann2021a}

% TODO: correlation models


\parencite{KrausWeyermann2021a}
SRM=L+0,04662*(L) hoch 2
L= 0,808*SRM-0,0083 für Helle
Kongressmaische, standartisierte Maischeprogramm
EBC 1996 EBC Skala 25-fache log. 1EBC = srm * 1,97,
3,6-25,3 zuverlässig. 
Dunkle Flüssigkeit absorbiert Viel. Keine zuverlässige Messung ab 80-100 EBC
80 EBC Wahrnehmungsschlelle. Dunkler als 80 wird nicht mehr vom Auge wahrgenommen.
Mit wasser verdünnen. 

\parencite{KrausWeyermann2021b}
Stark divergierende Farbwerte für gleiche Schüttungen.
Gleiche Farbwerte für unterschiedliche Schüttungen die 
unterschiedliche Schüttungen verursachen.
1L = 2,65 EBC
EBC=0,377L
bei mehreren Malzen proportionale Anteil der Schüttung.
Malt Color Unit MCU=Malzfarbe[°L]*Schüttungsanteil[lb]/Volumen Anstellwürze [gal]
lb gal 
1MCU = 1 SRM (bis 6 SRM)
cEBC=Malzfarbe[EBC]*Gewicht[kg]/Gesamtgewicht Maische * Stw/10 
Je dunkler, desto mehr faktoren sind zu berechnen
sobald bierfarbe 10 SRM überschreitet wird lineare MCU unbrauchbar

CEBC = G/10+BT*D+CC
CW= G/10+CC

Krüger hat Zufärbungskorrekturen
1-2 EBC pro Stunde kochen
Morey hat sich zum Standard entwickelt.
Mosher Kommerzielle Bire mit bekannten Schüttungen und Farbwerten. über 10 SRM nich sehr zuverlässig.
Daniels; versuche mit Hobbysuden

Morey: abnehmende Genauigkeit der Lineare für dunkel logarithmisch der Lichtabsoption
anwendbar bis 50 SRM

\section*{Homebrew/Malt Color Unit (HCU/MCU)}

\parencite[11]{Noonan1996}
HCU color units are a summary measurement of
wort color that is equal to the sum of the SRM/°Lovibond
of all the malts used to make a wort. HCU units and SRM
wort color are comparable up to about 10 °SRM; after that,
corrections must be made.


„\href{https://bieranalyse.de}{Bieranalyse Fuchs}“

%In \autoref{table:parameters} sind die vom \citeauthor{BJCP2015} und die von
%\citeauthor{Strottner1999} genannten brautechnischen Parameter erfasst.
%Die Angaben der Brewers Association Stil-Richtlinien sind annähernd
%deckungsgleich \parencites[24]{BA2021}. 

\parencite{Colby2000}
The HCU equation is similar to the SGP and IBU calculations, except it is missing a utilization or efciency term.
Essentially, it is only a measure of how much color is put into the beer from the grains. How that color is
changed by the brewing process is ignored.
Ignoring the effect of the brewing process on beer color is unfortunate, since many brewing processes affect
wort and beer color. Some factors increase color. These include: high hop rates, longer boil times, ner grain
crushes and oxidation from wort abuse or aging. Fermentation, in contrast, reduces the color of the wort.
The equation could be improved if a factor was added to account for changes in color due to the brewing
process. In fact, two factors would be needed. Some color changes are modications of the grain color added to
the beer. Loss of color during fermentation, for example, depends on the amount of color initially in the wort.
Other color changes do not depend on the initial beer color. Some of the wort darkening during the boil is due
to the sugars caramelizing. In solution, the sugars are colorless. The color that sugar caramelization contributes
does not come from the roasting of the grain.
So, an equation that took both of these factors into consideration would need two new variables. The resulting
equation would be:

NHCU=[W(lbs) * CR(degrees L)*TF + AF]/V(gallon)
NHBU is New Homebrew Color Units. The two new factors are TF and AF. TF is the transforming factor, the
extent that the grain color is modidfied. AF is the additive factor, the amount of color that appears during the
brewing process or aging that is not dependent on roasted grain-derived color.

\section*{Modelle zur Korrelation zwischen MCU und Bierfarbe}

Farbvorhersage

\subsection*{Mosher (1994)}

\parencite{Morey}:
\begin{equation}
\SRM = 0,3 \cdot \MCU + 4,7
\label{eq:mcumosher}
\end{equation}

Homebrew Color Unit (HCU), digramm
\parencite[34]{Mosher1994}

Metric Malt Color Units = l*kg/l
\parencite[258]{Mosher2015}

\subsection*{Daniels (1995)}

In seinem Buch „\citetitle{Daniels1996}“ hat Daniels die Korrelation zwischen SRM und MCU in \autoref{table:mcudaniels} veröffentlicht. \textcite[10]{Holle2010} verwendet die gleiche Tabelle für Farbvorhersagen. Die darauf basierende lineare \autoref{eq:mcudaniels} ist darin nicht enthalten, wird von \textcite{Morey} aber ebenfalls Daniels zugeschrieben. Eine in die Tabellenwerte eingepasste alternative Exponentationfunktion \autoref{eq:mcudanielsdruey} stammt von \textcite{Druey1998}.

\begin{table}[H]
\centering
\begin{tabular}{rrr}
\toprule
\multicolumn{1}{c}{\textbf{SRM}} & \multicolumn{1}{c}{\textbf{EBC}} & \multicolumn{1}{c}{\textbf{MCU}} \\
\midrule
1–10 & 2–20 & 1–10 \\
11–20 & 22–39 & 8–12 \\
21–30 & 41–59 & 11–15 \\
31–40 & 61–79 & 14–17 \\
41–50 & 81–98 & 17–20 \\
50–85 & 98–167 & 20–30 \\
>85 & >167 & >30 \\
\bottomrule
\end{tabular}
\caption{Korrelation zwischen SRM und MCU \parencite[61]{Daniels1996}}
\label{table:mcudaniels}
\end{table}

\begin{equation}
\SRM = 0,2 \cdot \MCU + 8,4
\label{eq:mcudaniels}
\end{equation}

\begin{equation}
\SRM = 1,73 \cdot \MCU^{0,64} - 0,267
\label{eq:mcudanielsdruey}
\end{equation}

\subsection*{Noonan (1996)}

\begin{table}[H]
\centering
\begin{tabular}{rrr}
\toprule
\multicolumn{1}{c}{\textbf{SRM}} & \multicolumn{1}{c}{\textbf{EBC}} & \multicolumn{1}{c}{\textbf{MCU}} \\
\midrule
1–10 & 2–20 & 1–10 \\
10,5 & 21 & 10,8 \\
11   & 22 & 11,6 \\
11,5 & 23 & 12,4 \\
12   & 24 & 13,3 \\
12,5 & 25 & 14,1 \\
13   & 26 & 14,9 \\
13,5 & 27 & 17,7 \\
14   & 28 & 18,6 \\
14,5 & 29 & 20,5 \\
15   & 30 & 22,4 \\
15,5 & 31 & 24,3 \\
16   & 32 & 26,2 \\
16,5 & 33 & 28,1 \\
17   & 33 & 30 \\
17,5 & 34 & 32,9 \\
18   & 35 & 35,8 \\
18,5 & 36 & 38,8 \\
19   & 37 & 41,9 \\
19,5 & 38 & 45 \\
20   & 39 & 47,8 \\
\bottomrule
\end{tabular}
\caption{Korrelation zwischen SRM und MCU \parencite[206]{Noonan1996}}
\label{table:mcunoonan}
\end{table}

\parencite{Druey1998}:
Kyle Druey
\begin{equation}
\SRM = 15,03 \cdot \MCU^{0,27} - 15,53
\label{eq:mcunoonandruey}
\end{equation}

\subsection*{Morey (1996)}

Heimbrauer Daniel Morey war unzufrieden mit der Farbvorhersage, die er in seiner Brausoftware auf Basis der MCU ohne zusätzliche Korrelation integriert hatte. Diese berechnete seiner Meinung nach zu hohe SRM Werte. Im mittlerweile nicht mehr erhältlichem Brewing Techniques Magazin erfuhr er dann 1995 von der Mosher und Daniels Modell. Zunächst wollte er alle Korrelationen in ein Berechnungsmodell zusammenführen, entschied sich aber dann dazu, die ihm vorliegenden Daten so zu modifizieren, dass es möglich war eine seinen Annahmen entsprechende Kurve einzupassen. \parencite{Smith2010}

Das von Morey veröffentlichte Modell beruht auf folgenden Annahmen: Bis zu einem Betrag von 10 entspricht die MCU ungefähr dem zu erwartenden SRM Messwert. Ab 10 MCU gibt das Daniels Modell eine bessere Korrelation, über 37 MCU hingegen das Mosher Modell. Das menschliche Auge kann keinen Farbunterschiede über 40 SRM mehr feststellen. Die \autoref{eq:mcumorey} wurde daher nur bis 50 SRM ausgelegt.
\parencite{Morey}

\begin{equation}
\SRM = 1,49 \cdot \MCU^{0,69}
\label{eq:mcumorey}
\end{equation}

\subsection*{Krüger (2019)}



\subsection*{Hanghofer (2019)}

\parencite[78]{Hanghofer2019}



\subsection*{Weyermann (2021)}

\parencite{KrausWeyermann2021b}

\begin{table}[H]
\centering
\begin{tabular}{rr}
\toprule
\multicolumn{1}{c}{\textbf{Stammwürze}} & \multicolumn{1}{c}{\textbf{CC [EBC]}} \\
\midrule
0–7 & 0 \\
7,1–10 & 3 \\
10,1–15 & 5 \\
15,1–20 & 7 \\
>20 & 0 \\
\bottomrule
\end{tabular}
\caption{TODO \parencite{KrausWeyermann2021b}}
\label{table:ccweyermann}
\end{table}

\subsection*{Modellvergleich}

\begin{figure}[h]
\centering
\includegraphics[width=14cm]{graph_srm.pdf}
\caption{TODO (Ascher, 2022)}
\label{fig:mcucompare}
\end{figure}


\parencite{KrausWeyermann2021b}
20 Biere in Weyermann Braumanufaktur. Gemessen und Labo-Messung mit Formel verglichen.
wobei 1 Bier (Export hell) als ausreißwer entfernt wurde.
keine 100 prozen überinstimmung mit labormessung.
Weyermann und Krüger beste Ergebnis.

\begin{table}[H]
\centering
\begin{tabular}{lr}
\toprule
\multicolumn{1}{c}{\textbf{Modell}} & \multicolumn{1}{c}{\textbf{Mittlere Abw. [\%]}} \\
\midrule
Daniels-Linear & -14,4 \\
Krüger & 6,9 \\
Morey & -51,8 \\
Mosher-Linear & -57,2 \\
Weyermann & -2,1 \\
\bottomrule
\end{tabular}
\caption{TODO \parencite{KrausWeyermann2021b}}
\label{table:modelcompare}
\end{table}

Ausschlagvolumen: 275 Liter
Böhmisches Pilsner
11,7 °P
Tennenmalz
42,8 kg (3,75 EBC); Carapils® 2,3 kg
(4,5 EBC); Sauermalz 0,9 kg (6 EBC);
Carabohemian® 0,5 kg(195 EBC)

12,3 °P Böhmisches Dunkles
Böhmisches Pilsner Tennenmalz 38,9
kg (3,75 EBC); Carapils 2,4 kg (4,5 EBC);
Carabohemian 2,4 kg (195 EBC); Sauer-
malz 0,9 kg (6 EBC); Carafa® Typ III, 1,9
kg (1400 EBC)

\begin{table}[H]
\centering
\begin{tabular}{lrr}
\toprule
\multicolumn{1}{c}{\textbf{Modell}} & \multicolumn{1}{c}{\textbf{Pilsner [EBC]}} & \multicolumn{1}{c}{\textbf{Dunkles [EBC]}} \\
\midrule
Messwert nach EBC & 10 & 82 \\
Daniels-Druey & 7 & 34 \\
Daniels-Linear & 18 & 31 \\
Hanghofer & & \\
Krüger & 11 & 91 \\
Morey & 6 & 36 \\
Mosher-Linear & 11 & 32 \\
Noonan-Druey & 10 & 48 \\
Weyermann & 12 & 92 \\
\bottomrule
\end{tabular}
\caption{TODO \textcite{KrausWeyermann2021c}}
\label{table:modelcomparepilsner}
\end{table}

\section*{Zusammenfassung}


\parencite{Bruecklmeier2018}
\parencite{Holle2010}

\parencite{Smith2008}
\parencite{Bies2010}
\parencite{Tucker2017}

\parencite{KrausWeyermann2021a}
\parencite{KrausWeyermann2021b}
\parencite{KrausWeyermann2021c}

\parencite{Lange2016}
\parencite{Caro2019}
\parencite{Daniels2012}

Die relevanten Eckpunkte dieses Artikels sind:

\begin{itemize}
\item TODO
\end{itemize}

\printbibliography[title=Quellen]

\end{document}